\begin{frame}{모델은 파이프라인의 한 조각일 뿐}
  1.2절은 ``모델을 어떻게 정의하는가''였지만, 실전에서 모델 학습은
  전체 작업의 \textbf{작은 일부}에 불과하다. 실제 데이터 프로젝트의
  순환:
  \begin{enumerate}
    \item \kb{수집} --- 어디서 데이터를 구할 것인가. Week 8, 16 팀
      프로젝트에서는 Kaggle, UCI ML Repository, Hugging Face Datasets 등
      공개 저장소를 주로 사용. 실무에선 ``이 데이터를 이 목적으로 써도
      되는가''(라이선스, 개인정보) 확인도 포함 --- 캐글 대회 데이터라도
      \textbf{상업적 재사용이 금지}된 경우가 흔함.
    \item \kb{정제} --- \textbf{결측치} · \textbf{중복} ·
      \textbf{이상치}를 처리 (현실 데이터는 거의 항상 지저분).
    \item \kb{탐색적 데이터 분석(EDA)} --- 모델을 만들기 전에
      분포·상관관계·클래스 비율을 \textbf{눈으로} 확인. 이 단계의 발견이
      어떤 모델을 고를지까지 결정 (예: 99:1 불균형 $\to$
      정확도 대신 PR-AUC, Week 2.3).
    \item \kb{모델링} --- 1.2절에서 정식화한 문제에 맞는 $h$ 를 고르고
      학습. 이번 학기 대부분의 장.
    \item \kb{평가} --- \textbf{학습에 쓰지 않은} 데이터로 성능 검증
      (엄밀한 원리는 Week 6). 나쁘면 4단계(다른 모델?)로만 돌아가지
      않고 \textbf{1$\sim$3단계로 되돌아가야} 할 수도 ---
      ``모델이 나쁜 것''이 아니라 ``데이터가 애초에 그 신호를 담고 있지
      않은 것''일 수 있기 때문.
  \end{enumerate}
\end{frame}
